\section{Background}
\label{sec:background}

\subsection{Medical Imaging Modalities and Dermoscopy Specifics}
Clinical imaging spans distinct physical acquisition mechanisms, each imposing specific preprocessing and invariance constraints. Projection radiography (chest X-rays) measures cumulative X-ray attenuation in grayscale intensities where gravitational anatomy forbids vertical flipping and contrast stretching (such as CLAHE) compensates for exposure variability. Computed tomography (CT) reconstructs calibrated Hounsfield Units requiring tissue-specific windowing (e.g., lung vs.\ soft-tissue windows), whereas magnetic resonance imaging (MRI) exhibits non-standardized signal scales sensitive to scanner magnetic field strength and coil bias. Histopathology captures gigapixel H\&E-stained whole-slide tiles whose color distribution varies with laboratory staining protocols~\cite{Yang2023medmnist}, while dense pixel-level lesion delineation across modalities frequently builds on nested convolutional architectures such as UNet++~\cite{Zhou2020unet}.

Dermoscopy (epiluminescence microscopy) occupies a distinct niche in visible-light clinical photography. A handheld dermatoscope combines $10\times$ optical magnification with cross-polarized illumination or immersion fluid to suppress specular reflection from the stratum corneum, exposing subsurface pigmented networks, globules, blue-white veils, and vascular structures critical to diagnosing melanoma, basal cell carcinoma, and actinic keratoses~\cite{Tschandl2018dataset}. Unlike chest radiographs, dermatoscopic fields of view possess approximate rotational and reflectional symmetry (the dihedral group $D_4$), making $90^\circ$ rotations and horizontal or vertical flips clinically valid augmentations. Conversely, chromatic cues—particularly brown-to-black melanin variation, erythema, and violaceous vascular lacunae—carry primary diagnostic signal under the clinical ABCDE and 7-point checklists; aggressive hue jittering or channel permutation therefore violates clinical validity. Three modality-specific artifacts regularly corrupt dermatoscopic images and induce shortcut learning: black circular optical vignettes at image corners, dark terminal hair strands occluding lesion borders, and illuminant color casts across different dermatoscope models and clinical centers~\cite{Lee1997dullrazor,Finlayson2004shades}.

\subsection{Clinical Workflows and Decision Context}
Dermatological diagnosis operates across a tiered care pathway. In primary care and community clinics, general practitioners encounter high volumes of benign pigmented lesions alongside uncommon high-risk malignancies. The clinical objective at this first point of contact is high-sensitivity triage: missing an invasive melanoma (a false negative) carries severe mortality risk, whereas referring an ambiguous benign lesion to dermatology incurs manageable downstream cost~\cite{Tschandl2019comparison}. Within specialist dermatology clinics, clinicians integrate dermoscopy with patient history (age, lesion evolution, personal or family history of melanoma, sun-exposure site, and comparison against the patient's other nevi) to decide whether to perform excisional biopsy, initiate digital dermoscopic follow-up, or reassure the patient.

\begin{figure}[htbp]
\centering
\resizebox{\textwidth}{!}{%
\begin{tikzpicture}[
  node distance=0.55cm and 0.65cm,
  stepbox/.style={rectangle, rounded corners=3pt, draw=OkabeBlue, fill=BoxBg, thick, text width=2.55cm, align=center, minimum height=1.25cm, font=\footnotesize},
  gate/.style={diamond, aspect=2.1, draw=OkabeOrange, fill=orange!8, thick, align=center, font=\footnotesize, inner sep=1pt},
  clinic/.style={rectangle, rounded corners=3pt, draw=OkabeGreen, fill=teal!7, thick, text width=2.7cm, align=center, minimum height=1.25cm, font=\footnotesize},
  arrow/.style={-{Stealth[length=2.2mm]}, thick, color=OkabeGray}
]
\node[stepbox] (acq) {\textbf{1. Dermoscopy Acquisition}\\ Polarized VL-DICOM / JPEG};
\node[stepbox, right=of acq] (phi) {\textbf{2. PHI Audit \& Preproc}\\ pydicom Allowlist\\ DullRazor + SoG};
\node[stepbox, right=of phi] (inf) {\textbf{3. Calibrated Inference}\\ $P(\text{mal})$ + Temp. Scale\\ + Grad-CAM / Rollout};
\node[gate, right=0.75cm of inf] (thr) {\textbf{$P(\text{mal}) \ge \tau_{\text{scr}}$?}};
\node[clinic, above right=0.45cm and 0.7cm of thr] (high) {\textbf{Urgent Dermatology Referral}\\ Biopsy / Excision Review};
\node[clinic, below right=0.45cm and 0.7cm of thr] (low) {\textbf{Routine Follow-up}\\ Clinician ABCDE Override Allowed};

\draw[arrow] (acq) -- (phi);
\draw[arrow] (phi) -- (inf);
\draw[arrow] (inf) -- (thr);
\draw[arrow] (thr.north) |- node[pos=0.72, above, font=\scriptsize]{Yes (Sens $\ge 95\%$)} (high.west);
\draw[arrow] (thr.south) |- node[pos=0.72, below, font=\scriptsize]{No (High NPV)} (low.west);
\end{tikzpicture}
}
\caption{\textbf{Proposed human-in-the-loop primary-care teledermatology triage workflow.} Calibrated malignancy probabilities operate at a validation-fixed high-sensitivity threshold ($\tau_{\text{scr}}$) paired with Grad-CAM visual explanations, ensuring dermatologists and primary-care physicians retain final diagnostic authority and override capability.}
\label{fig:clinical_workflow}
\end{figure}


As illustrated in Figure~\ref{fig:clinical_workflow}, an AI system for skin lesion assessment cannot operate as an autonomous diagnostic oracle. Instead, it fits into teledermatology and primary-care triage as a calibrated second reader: lesions exceeding a high-sensitivity screening threshold trigger prioritized specialist evaluation, while visual saliency overlays allow clinicians to inspect whether the prediction focuses on true lesion morphology or an extraneous border artifact. Prospective human-computer interaction trials on ISIC~2018 lesions showed that accurate multiclass AI decision support improved clinician diagnostic accuracy by 13~percentage points—particularly among less experienced practitioners—yet faulty AI predictions degraded clinician accuracy even among experts through automation bias~\cite{Tschandl2020human}.

\subsection{Regulatory, Privacy, and Quality Standards Framework}
Clinical AI software intended to inform diagnosis, screening, or treatment of disease is regulated in the United States as Software as a Medical Device (SaMD) by the U.S.\ Food and Drug Administration (FDA)~\cite{Food2021artificial,Gerke2020need}. Three primary premarket pathways govern authorization depending on risk class and availability of legally marketed predicates: (i)~\emph{510(k) clearance} for moderate-risk Class~II devices demonstrating substantial equivalence to an existing predicate device; (ii)~\emph{De~Novo classification} for novel low-to-moderate-risk devices lacking a predicate—frequently used for first-in-kind AI triage or dermatology CADx tools; and (iii)~\emph{Premarket Approval (PMA)} for high-risk Class~III autonomous or life-critical diagnostics requiring extensive prospective clinical trials~\cite{Zhu2022landscape}. Laboratory-Developed Tests (LDTs) historically operated under CLIA enforcement discretion inside single clinical laboratories, though FDA oversight over software-driven LDTs has tightened considerably and does not cover commercially distributed imaging software.

Crucially, conventional SaMD approvals lock model weights at clearance. Because medical imaging distributions evolve with scanner replacements, software upgrades, and demographic shifts~\cite{Finlayson2021clinician}, the FDA AI/ML SaMD Action Plan~\cite{Food2021artificial} introduced Good Machine Learning Practice (GMLP) principles and Predetermined Change Control Plans (PCCPs), allowing manufacturers to pre-specify retraining protocols, validation metrics, and guardrails for post-market updates without full re-submission. Complementary international standards include ISO~13485 (medical device quality management systems), IEC~62304 (medical device software lifecycle), DICOM PS3.15 (Visible Light Photography SOP Class \texttt{1.2.840.10008.5.1.4.1.1.77.1.4} and de-identification confidentiality profiles), and HIPAA Safe Harbor rules (45~CFR \S164.514(b)) mandating removal of 18~patient identifiers alongside protection against re-identification via burned-in pixel text or facial features.
